\begin{afterword}
\summaryhead

The post distinguishes two styles of thought. ``Toolbox'' thinkers keep many methods and choose among them by context. ``Law'' thinkers hold that some truths govern every case, including ideals that no practical method reaches.

It starts from a Twitter exchange in which David Chapman defined rationalism as the claim that there is one correct way to think, and Julia Galef replied that rationalists hold a normative ideal rather than a single rule. The author then stages a dialogue between the two styles about statistics.

A long maze analogy follows. Every maze has a shortest path. A walker may do better to follow a simple rule, and knowing that a shortest path exists shows where improvement is possible. The author grants that law thinkers can forget context, but argues that both styles are tools to be used as context demands. The post closes by suggesting that laws such as the triangle inequality, Carnot's limit and Bayesian updating are best understood as plain facts rather than ideals.

\responsehead

The post is a reply to a named critic. It says openly that it will ``badly stereotype this conversation'' as the author has seen it many times, and it builds its toolbox character from that composite rather than from Chapman. Chapman's tweets, as the post quotes them, already describe the Law position: rationalisms that ``insist there must be one, but that it is not currently knowable''. When the distinction between an ideal (normative) and a working recipe (prescriptive) is put to Chapman, Chapman calls it ``useful and maybe key'' and names the real disagreement: whether an ideal that ``can't be achieved, even in principle, is a good thing.'' The post's diagnosis is that toolbox thinkers cannot see what an ideal is for except as a recipe. Its only evidence shows a toolbox thinker who saw the point at once and asked a harder question.

That question is never answered, because the central analogy is built so that it cannot arise. A maze comes with an agreed measure, distance, and its shortest path can be walked by anyone with a map. What the measure for reasoning is, and whether its ideal can be reached even in principle, are the two things in dispute. The maze assumes both answers. The post's positive claim, that thinking about ideals gets you to some destinations faster, is asserted without one example.

The opposition is written by the author, and written to lose. Its objections are phrased as universals (``surely'', ``inevitably'', ``never''), so that ``that happens some of the time'' can dispose of them. The one real case in the post, the mortgage-security models blamed for the 2008 crisis, is a failure of Law thinking, is put in the opponent's mouth, and is not discussed.

The post breaks its own rule. It concedes that Law thinkers forget the premise between ``shortest'' and ``best''. Then it states that no update can beat Bayes in expected score, leaving out the premise that the prior and likelihoods are right, the kind of premise the mortgage models got wrong.

The framing is not new either. The distinction between descriptive, normative and prescriptive models dates from the 1980s in decision research, and the post credits no one. Its resolution, that both styles are tools to be used by context, is the one it had itself attributed to the ``Toolbox-inclined''. Throughout, the post ranks its readers: ``deeper'' and ``more advanced'' understanding lies on the Law side, and toolbox thinkers get a closing remedy for their confusion.

In short: a reply to a critic that stereotypes the exchange by its own admission, argues from a maze constructed so that the ideal Chapman doubted exists, and ends at the answer it had assigned to the toolbox side, while ranking itself above it.
\end{afterword}
